\documentclass[11pt]{article}
\usepackage[T1]{fontenc}
\usepackage[utf8]{inputenc}
\usepackage[english]{babel}
\usepackage{geometry}
\geometry{margin=2.5cm}
\usepackage{booktabs}
\usepackage{array}
\usepackage{longtable}
\usepackage{graphicx}
\graphicspath{{figures/}}
\usepackage{xr}
\externaldocument{manuscript}
\usepackage{amsmath}
\usepackage{siunitx}
\usepackage{hyperref}
\hypersetup{colorlinks=true,linkcolor=blue,urlcolor=blue,citecolor=blue}
\usepackage[strings]{underscore}
\usepackage{fancyhdr}

\setlength{\headheight}{14pt}
\pagestyle{fancy}
\fancyhf{}
\fancyhead[L]{Supporting Information --- $\Delta$(ICOBI\textsubscript{antibonding})}
\fancyhead[R]{\thepage}
\renewcommand{\headrulewidth}{0.4pt}

\renewcommand{\thetable}{S\arabic{table}}
\renewcommand{\thefigure}{S\arabic{figure}}
\renewcommand{\thesection}{S\arabic{section}}
\renewcommand{\theequation}{S\arabic{equation}}

\title{\bfseries Supporting Information\\[4pt]
\large $\Delta$(ICOBI\textsubscript{antibonding}): A Fast Bond-Order-Based
Descriptor for the Kinetic Viability of Metastable Crystalline Compounds}
\author{Gilles Frapper \\ IC2MP, UMR 7285 CNRS, Universit\'e de Poitiers}
\date{}

\begin{document}
\maketitle
\tableofcontents
\newpage

\section{Dataset Composition and Selection Criteria}
\label{sec:si-dataset}

Candidate structures were drawn from the Materials Project database by
querying unary and binary chemical systems, restricted to non-magnetic
structures (absolute total magnetization $\leq 0.01~\mu_\text{B}$ per
formula unit, computed directly from the stored magnetic moments) and
excluding $f$-block elements. Three complementary selection tracks were
combined into the final 597-structure dataset:

\begin{itemize}
  \item \textbf{Track 1} (186 compounds): $\leq 20$ atoms/cell, balanced
  across stability classes (60 per family), with a cap of two entries per
  chemical system to maximize elemental diversity over redundant polymorphs.
  \item \textbf{Track 2} (163 compounds/elements): $\leq 40$ atoms/cell,
  targeting additional elemental references, off-equilibrium polymorphs
  (high-pressure phases, allotropes), and alkali/alkaline-earth binaries
  against N, O, F, P, S, and Cl, chosen to build genuine same-composition
  polymorph groups and probe descriptor behavior far from equilibrium.
  \item \textbf{Track 3} (45 binary compounds): an adversarial robustness
  set pairing an experimentally known compound sitting farthest above the
  convex hull of its chemical system with one to three theoretical
  polymorphs of the same formula lying even farther above the hull.
\end{itemize}

Table~\ref{tab:si-composition} summarizes the dataset by element count and
Table~\ref{tab:si-stability} by stability class. One trivial-cell reference
element yielded no usable LOBSTER step and is excluded from every
ICOHP/ICOBI-derived statistic (Table~\ref{tab:si-composition}, footnote).

\begin{table}[htbp]
\centering
\caption{Dataset composition by element count.}
\label{tab:si-composition}
\begin{tabular}{@{}lr@{}}
\toprule
Composition & $n$ \\
\midrule
Single-element references & 69 \\
Multi-element compounds   & 522 \\
Additional structures without compound metadata\textsuperscript{a} & 6 \\
\midrule
\textbf{Total} & \textbf{597} \\
\bottomrule
\end{tabular}
\vspace{0.15cm}
\\{\footnotesize \textsuperscript{a} isolated gas dimers (H$_2$, N$_2$,
O$_2$, F$_2$, Cl$_2$) and a CaO sphalerite-type reference, used only for
internal quality-control checks (Section~\ref{sec:si-qc}), outside the main
dataset.}
\end{table}

\begin{table}[htbp]
\centering
\caption{Stability class, by $E_\text{hull}$ and experimental/theoretical
character.}
\label{tab:si-stability}
\begin{tabular}{@{}lrl@{}}
\toprule
Class & $n$ & $E_\text{hull}$ range (eV/atom) \\
\midrule
Stable ($E_\text{hull}=0$) & 160 & 0.0000 \\
Experimental metastable    & 183 & 0.0000347 -- 2.1080 \\
Theoretical metastable     & 242 & 0.0000843 -- 3.1688 \\
Unknown\textsuperscript{a} & 12 & --- \\
\midrule
\textbf{Total} & \textbf{597} & \\
\bottomrule
\end{tabular}
\vspace{0.15cm}
\\{\footnotesize \textsuperscript{a} entries lacking a usable
$E_\text{hull}$ value or predating the \texttt{theoretical} flag used to
build later batches of this dataset.}
\end{table}

\section{VASP Calculation Parameters}
\label{sec:si-incar}

Table~\ref{tab:si-incar} lists the common INCAR parameters applied
identically to every structure in the dataset (Section~2.3 of the main
text). Table~\ref{tab:si-potcar} lists the PAW-PBE potential used for each
element (\texttt{pymatgen}'s \texttt{MPRelaxSet} library, with plain
\texttt{W} substituted for \texttt{W\_sv}/\texttt{W\_pv} after the latter was
found to be missing its \texttt{LEXCH} field in the local potential mirror;
Section~\ref{sec:si-qc}).

\begin{table}[htbp]
\centering
\caption{Common VASP INCAR parameters.}
\label{tab:si-incar}
\begin{tabular}{@{}l l p{7.5cm}@{}}
\toprule
Parameter & Value & Description \\
\midrule
\texttt{PREC} & Accurate & Overall precision (charge grid, projectors) \\
\texttt{ENCUT} & 600.0 & Plane-wave cutoff energy (eV) \\
\texttt{EDIFF} & 1e-06 & Electronic convergence criterion (eV) \\
\texttt{IBRION} & -1 & Ionic dynamics type ($-1$ = static, no relaxation) \\
\texttt{NSW} & 0 & Number of ionic steps (0 = static run on the MP structure) \\
\texttt{ISMEAR} & 0 & Smearing type ($0$ = Gaussian) \\
\texttt{SIGMA} & 0.05 & Smearing width (eV) \\
\texttt{LASPH} & True & Non-spherical corrections to the PAW potential \\
\texttt{ISYM} & -1 & Symmetry ($-1$ = disabled, required by LOBSTER) \\
\texttt{LWAVE} & True & Write WAVECAR (required by LOBSTER) \\
\texttt{LCHARG} & True & Write CHGCAR \\
\texttt{NPAR} & 4 & Band-level parallelization \\
\texttt{KPAR} & 2 & K-point-level parallelization \\
\bottomrule
\end{tabular}
\end{table}

{\small
\input{si_potcar_body.tex}
}

\section{$k$-Point Grid and NBANDS per Compound}
\label{sec:si-kpoints}

Table~\ref{tab:si-kpoints} lists, for the 186 compounds of Track~1
(Section~\ref{sec:si-dataset}), the automatic $k$-point grid
(\texttt{pymatgen} automatic density, \texttt{kppvol=100}; G/M = $\Gamma$-
or Monkhorst--Pack-centered) and the number of bands (\texttt{NBANDS},
derived from the number of local orbitals required by the LOBSTER basis for
that compound). The same generation rule was applied to every other track
of the dataset; the complete machine-readable listing for all 597
structures is available in the project's data repository (Data
Availability, main text).

{\small
\input{si_kpoints_body.tex}
}

\section{Elemental Reference Structures}
\label{sec:si-elements}

Table~\ref{tab:si-elements} lists the 62 elemental reference structures used
as the products side of every case-1 (decomposition-into-elements)
reaction, together with their raw and near-frontier antibonding ICOHP
values (main text, Section~2.6).

{\small
\input{si_elements_body.tex}
}

\section{Bonding-Type Classification: Full Counts}
\label{sec:si-bondtype}

Table~\ref{tab:si-bondtype-full} reports, for the two independent bonding
classifications discussed in the main text (Section~3.6), the full
cross-tabulation over all 588 compounds and elements with usable
ICOHP/ICOBI data.

\begin{table}[htbp]
\centering
\caption{Bonding-type classification: composition-based heuristic vs.
ICOHP/ICOBI-based classification ($n=588$).}
\label{tab:si-bondtype-full}
\begin{tabular}{@{}lrrrrr@{}}
\toprule
Heuristic $\backslash$ ICOHP/ICOBI & metallic & ionic & covalent & mixed & unclassified \\
\midrule
metallic     & 193 & 0  & 0  & 0  & 0 \\
ionic        & 0   & 48 & 2  & 3  & 0 \\
covalent     & 0   & 14 & 35 & 16 & 0 \\
unclassified & 179 & 36 & 19 & 43 & 0 \\
\midrule
\textbf{Total} & \textbf{372} & \textbf{98} & \textbf{56} & \textbf{62} & \textbf{0} \\
\bottomrule
\end{tabular}
\vspace{0.15cm}
\\{\footnotesize Concordance on the subset where the composition-based
heuristic renders a verdict: 276/311 = 88.7\%. The remaining 277 compounds
the heuristic leaves ``unclassified'' (mostly alkali/alkaline-earth $\times$
\{N,P,S\} pairs outside a fixed anion list) are resolved by the
ICOHP/ICOBI-based classification into 179 metallic, 36 ionic, 19 covalent,
and 43 mixed (Zintl-type).}
\end{table}

The mixed (Zintl-type) category identifies compounds where a
\emph{homoatomic} pair above the covalent ICOBI threshold (e.g. N--N in an
azide, P--P in a polyphosphide, S--S in a polysulfide) coexists with a
markedly weaker \emph{heteroatomic} pair linking a distinct cation, the
structural signature of Zintl--Klemm chemistry. The complete per-compound
classification for all 588 entries, listing formula, Materials Project
identifier, space group, dominant-pair ICOBI/ICOHP, and the normalized
ICOHP/ICOBI antibonding population, is reproduced by bonding-type category
in Sections~\ref{sec:si-master-metallic}--\ref{sec:si-master-unclassified}.

\subsection{Metallic Compounds and Elements}
\label{sec:si-master-metallic}
{\footnotesize
\begin{longtable}{@{}lllrrrr@{}}
\caption{Metallic compounds and elements (\texttt{is\_metal=True}, observed ICOBI 0.0000 to 2.9506) --- formula, mp-id, space group, dominant (first-shell) bond ICOBI, dominant bond ICOHP, normalized ICOHP and ICOBI antibonding populations ($\Delta E=1.0$~eV).}\\
\toprule
Formula & mp-id & Space group & ICOBI & ICOHP (eV) & ICOHP antibonding & ICOBI antibonding \\
\midrule
\endfirsthead
\toprule
Formula & mp-id & Space group & ICOBI & ICOHP (eV) & ICOHP antibonding & ICOBI antibonding \\
\midrule
\endhead
\bottomrule
\endfoot
\bottomrule
\multicolumn{7}{@{}p{0.95\linewidth}}{\footnotesize Black $^*$ = stable experimental compound (thermodynamic phase); red \textcolor{red}{$^*$} = metastable experimental compound; unmarked = theoretical. ICOBI/ICOHP: value of the dominant first-shell species pair (not a mean over every LOBSTER-reported neighbor, see \S\ref{sec:bondtype}).} \\
\endlastfoot
Ag\textcolor{red}{$^*$} & mp-124 & Fm-3m & 0.1545 & -1.2496 & 0.0000 & 0.0000 \\
Ag$_{3}$As & mp-1183193 & P6\_3/mmc & 0.2236 & -1.3102 & 0.0000 & 0.0000 \\
Ag$_{3}$Pt\textcolor{red}{$^*$} & mp-12065 & Pm-3m & 0.2113 & -1.5965 & 0.0033 & 0.0026 \\
Ag$_{3}$Sb\textcolor{red}{$^*$} & mp-2273 & Pmm2 & 0.2406 & -1.5221 & 0.0005 & 0.0000 \\
AgI\textcolor{red}{$^*$} & mp-22915 & Pm-3m & 0.2126 & -1.1434 & 0.0699 & 0.0000 \\
AgN\textcolor{red}{$^*$} & mp-1091417 & Cmce & 0.4624 & -3.4181 & 0.0088 & 0.0030 \\
AgN & mp-1009766 & F-43m & 0.4854 & -3.5762 & 0.0277 & 0.0173 \\
AgN & mp-1424734 & Ama2 & 0.8238 & -5.9446 & 0.0107 & 0.0025 \\
AgN & mp-1428312 & Pm-3m & 0.2177 & -1.4991 & 0.0083 & 0.0041 \\
AgPb$_{3}$ & mp-1183256 & I4/mmm & 0.2181 & -1.2748 & 0.0000 & 0.0000 \\
AgPd & mp-1183222 & P-6m2 & 0.1929 & -1.4704 & 0.0292 & 0.0199 \\
Al$^*$ & mp-134 & Fm-3m & 0.2270 & -1.4193 & 0.0000 & 0.0000 \\
Al$_{12}$Mo\textcolor{red}{$^*$} & mp-550 & Im-3 & 0.3238 & -1.5619 & 0.0000 & 0.0000 \\
Al$_{12}$W\textcolor{red}{$^*$} & mp-11227 & Im-3 & 0.5027 & -2.7492 & 0.0000 & 0.0000 \\
Al$_{2}$W\textcolor{red}{$^*$} & mp-12524 & P6\_422 & 0.4364 & -2.4880 & 0.0000 & 0.0000 \\
Al$_{3}$Ag & mp-1183130 & Pm-3m & 0.2444 & -1.4593 & 0.0000 & 0.0000 \\
Al$_{3}$As & mp-1183129 & Pm-3m & 0.2253 & -1.3902 & 0.0001 & 0.0000 \\
Al$_{3}$Os$_{2}$\textcolor{red}{$^*$} & mp-16521 & I4/mmm & 0.3022 & -1.6476 & 0.0037 & 0.0000 \\
Al$_{4}$Si & mp-1228120 & R-3m & 0.2614 & -2.0375 & 0.0087 & 0.0030 \\
AlAg & mp-1228901 & P-6m2 & 0.2436 & -1.7117 & 0.0000 & 0.0000 \\
AlAg$_{2}$ & mp-985805 & Cmcm & 0.2723 & -1.5388 & 0.0000 & 0.0000 \\
AlNi$^*$ & mp-1487 & Pm-3m & 0.2676 & -1.9848 & 0.0016 & 0.0023 \\
AlSb & mp-1023918 & Cmcm & 0.2647 & -1.2802 & 0.0000 & 0.0000 \\
As$^*$ & mp-158 & Cmce & 0.8848 & -3.9683 & 0.0103 & 0.0000 \\
As$_{3}$W$_{2}$$^*$ & mp-11696 & C2/m & 0.7072 & -4.3348 & 0.0000 & 0.0000 \\
AsPd$_{2}$\textcolor{red}{$^*$} & mp-1080831 & P-62m & 0.3410 & -2.4185 & 0.0077 & 0.0157 \\
AsRh\textcolor{red}{$^*$} & mp-22079 & Pnma & 0.3323 & -1.8639 & 0.0449 & 0.0158 \\
AsRh$_{2}$\textcolor{red}{$^*$} & mp-1102364 & Pnma & 0.2438 & -1.3778 & 0.1158 & 0.0189 \\
Au$^*$ & mp-81 & Fm-3m & 0.1962 & -1.3381 & 0.0000 & 0.0000 \\
B$_{3}$Ir$_{2}$ & mp-1228647 & C2/m & 0.4473 & -3.3600 & 0.0047 & 0.0011 \\
BPd$_{5}$ & mp-1228665 & I4/mmm & 0.4754 & -3.9868 & 0.0971 & 0.0181 \\
BPd$_{5}$ & mp-1228665 & I4/mmm & 0.4754 & -3.9924 & 0.0973 & 0.0181 \\
BW$^*$ & mp-7832 & I4\_1/amd & 0.4277 & -3.9556 & 0.0000 & 0.0000 \\
Ba$^*$ & mp-122 & Im-3m & 0.0002 & -0.0122 & 8.2400 & 2.9573 \\
Ba$_{3}$Ca & mp-1183342 & I4/mmm & 0.0000 & 0.0000 & -- & -- \\
Ba$_{3}$N$_{2}$ & mp-1013528 & Pm-3m & 0.0763 & -0.6612 & 0.1655 & 0.0007 \\
BaAu$_{2}$$^*$ & mp-30363 & P6/mmm & 0.6152 & -3.0712 & 0.0144 & 0.0034 \\
BaF$_{3}$ & mp-1183303 & I4/mmm & 0.0350 & -0.5277 & 0.2805 & 0.2537 \\
BaMg$_{3}$ & mp-1016218 & Pmm2 & 0.0023 & -0.0755 & 5.4796 & 1.1861 \\
BaN$_{2}$\textcolor{red}{$^*$} & mp-1102371 & Cc & 0.1605 & -0.4471 & 0.1039 & 0.0214 \\
BaSn$_{5}$\textcolor{red}{$^*$} & mp-571353 & P6/mmm & 0.2221 & -0.8507 & 0.0456 & 0.0042 \\
Be$^*$ & mp-87 & P6\_3/mmc & 0.2033 & -1.8975 & 0.0000 & 0.0000 \\
Be$_{12}$Mo$^*$ & mp-1475 & I4/mmm & 0.2043 & -1.9103 & 0.0000 & 0.0000 \\
Be$_{12}$Pd\textcolor{red}{$^*$} & mp-12666 & I4/mmm & 0.2078 & -1.9466 & 0.0000 & 0.0000 \\
Be$_{2}$Cu$^*$ & mp-2031 & Fd-3m & 0.3616 & -3.0863 & 0.0000 & 0.0001 \\
Be$_{2}$Nb$_{3}$$^*$ & mp-11272 & P4/mbm & 0.4200 & -3.5436 & 0.0000 & 0.0000 \\
Be$_{2}$V$^*$ & mp-11281 & P6\_3/mmc & 0.3658 & -2.7859 & 0.0000 & 0.0000 \\
BeAu$_{3}$ & mp-983539 & Pm-3m & 0.2230 & -1.9129 & 0.0000 & 0.0030 \\
BeCu\textcolor{red}{$^*$} & mp-2323 & Pm-3m & 0.1984 & -1.8359 & 0.0000 & 0.0000 \\
BeF$_{2}$ & mp-975644 & P6/mmm & 0.1548 & -1.6978 & 0.0863 & 0.0000 \\
Bi$_{2}$O$_{3}$\textcolor{red}{$^*$} & mp-558953 & P2\_1/c & 0.9517 & -6.7444 & 0.0328 & 0.0038 \\
Bi$_{2}$O$_{3}$ & mp-1383505 & P1 & 0.9506 & -7.1985 & 0.0609 & 0.0000 \\
C\textcolor{red}{$^*$} & mp-48 & P6\_3/mmc & 1.2195 & -11.5387 & 0.0000 & 0.0000 \\
Ca\textcolor{red}{$^*$} & mp-21 & Im-3m & 0.0025 & -0.0304 & 3.2206 & 3.3021 \\
Ca$_{3}$N$_{2}$ & mp-1013524 & Pm-3m & 0.1107 & -0.3623 & 0.0090 & 0.0000 \\
Ca$_{3}$Zn & mp-1039229 & I4/mmm & 0.1222 & -0.3874 & 0.2004 & 0.0086 \\
CaAl$_{4}$$^*$ & mp-1749 & I4/mmm & 0.4595 & -1.8701 & 0.0086 & 0.0002 \\
CaGa$_{2}$\textcolor{red}{$^*$} & mp-992 & P6/mmm & 0.9791 & -3.8887 & 0.0219 & 0.0000 \\
CaGa$_{4}$\textcolor{red}{$^*$} & mp-2461 & C2/m & 0.4641 & -2.1536 & 0.0084 & 0.0001 \\
CaGe$_{2}$ & mp-643542 & C2/m & 0.7379 & -3.1645 & 0.0337 & 0.0038 \\
CaN\textcolor{red}{$^*$} & mp-1058549 & Fm-3m & 0.1201 & -0.6934 & 0.0450 & 0.0656 \\
CaNi\textcolor{red}{$^*$} & mp-1059707 & Pm-3m & 0.0678 & -0.1910 & 0.5038 & 0.0500 \\
CaO & mp-1181903 & Fd-3m & 0.0725 & -0.5516 & 0.1865 & 0.0756 \\
CaSn$^*$ & mp-2450 & Cmcm & 0.7044 & -2.1965 & 0.0629 & 0.0276 \\
Cd$^*$ & mp-94 & P6\_3/mmc & 0.2423 & -1.6119 & 0.0000 & 0.0000 \\
CdIn$_{3}$ & mp-1183674 & Fm-3m & 0.2417 & -1.4735 & 0.0000 & 0.0000 \\
CdPd$^*$ & mp-1696 & P4/mmm & 0.2022 & -1.5904 & 0.0000 & 0.0004 \\
CdSb$_{3}$ & mp-1183739 & P6\_3/mmc & 0.2091 & -1.1803 & 0.0378 & 0.0058 \\
Co$^*$ & mp-102 & Fm-3m & 0.2680 & -1.9150 & 0.0189 & 0.0373 \\
CoB$^*$ & mp-20857 & Pnma & 0.4568 & -4.5928 & 0.0099 & 0.0000 \\
CoGe$^*$ & mp-21237 & C2/m & 0.4261 & -2.9354 & 0.0002 & 0.0022 \\
CoRe & mp-1225995 & P-6m2 & 0.3611 & -2.2350 & 0.0086 & 0.0082 \\
CoSe$_{2}$ & mp-1390196 & Fd-3m & 0.7048 & -4.3267 & 0.0348 & 0.0079 \\
CoSn$^*$ & mp-20536 & P6/mmm & 0.4900 & -2.7764 & 0.0035 & 0.0023 \\
Cr$^*$ & mp-90 & Im-3m & 0.2487 & -0.6245 & 0.0115 & 0.0049 \\
Cr$_{2}$C & mp-1226378 & P-3m1 & 0.5270 & -3.0803 & 0.0084 & 0.0000 \\
Cr$_{5}$Si$_{3}$\textcolor{red}{$^*$} & mp-7506 & I4/mcm & 0.4031 & -3.1831 & 0.0057 & 0.0003 \\
CrMo & mp-1226200 & Cmmm & 0.3600 & -1.3679 & 0.0081 & 0.0026 \\
CrRe$_{3}$ & mp-1183703 & Pm-3m & 0.3358 & -1.4730 & 0.0060 & 0.0000 \\
CrW & mp-1226162 & Cmmm & 0.5827 & -3.4455 & 0.0000 & 0.0005 \\
Cs\textcolor{red}{$^*$} & mp-1 & Im-3m & 0.0556 & -0.1230 & 0.1696 & 0.0693 \\
CsO$_{2}$\textcolor{red}{$^*$} & mp-1441 & I4/mmm & 1.0942 & -11.4627 & 0.1143 & 0.1436 \\
CsO$_{2}$ & mp-1096936 & Cmmm & 1.1298 & -11.0556 & 0.1197 & 0.0872 \\
CsRb & mp-1225895 & Cmmm & 0.0628 & -0.1628 & 0.1759 & 0.0961 \\
Cu$^*$ & mp-30 & Fm-3m & 0.0456 & -0.2488 & 0.1641 & 0.0946 \\
Cu$_{2}$S\textcolor{red}{$^*$} & mp-12087 & Fm-3m & 0.1495 & -0.9496 & 0.1576 & 0.0012 \\
Cu$_{2}$S\textcolor{red}{$^*$} & mp-553942 & P6\_3/mmc & 0.0830 & -0.4779 & 0.1504 & 0.0087 \\
Cu$_{3}$Ge\textcolor{red}{$^*$} & mp-19724 & Pmmn & 0.1369 & -0.8340 & 0.0963 & 0.0410 \\
Cu$_{3}$Te$_{2}$ & mp-1225676 & P4mm & 0.1480 & -0.8199 & 0.1597 & 0.0663 \\
Cu$_{7}$Te$_{4}$\textcolor{red}{$^*$} & mp-624307 & P3m1 & 0.1323 & -0.6945 & 0.1434 & 0.0510 \\
CuH & mp-1225705 & P-3m1 & 0.0446 & -0.1622 & 0.1097 & 0.0037 \\
CuNi & mp-1184069 & P-6m2 & 0.1198 & -0.4298 & 0.4478 & 0.0732 \\
Fe$^*$ & mp-13 & Im-3m & 0.1417 & -0.5370 & 0.0922 & 0.0794 \\
FeS$^*$ & mp-505531 & P4/nmm & 0.4967 & -2.5883 & 0.0736 & 0.0024 \\
FeTe$_{2}$\textcolor{red}{$^*$} & mp-1102002 & Pa-3 & 0.4279 & -1.7861 & 0.0401 & 0.0000 \\
Ga$^*$ & mp-142 & Cmce & 0.4308 & -2.6378 & 0.0000 & 0.0000 \\
Ga$_{2}$Au$_{7}$\textcolor{red}{$^*$} & mp-571598 & P-62m & 0.2592 & -1.8968 & 0.0000 & 0.0000 \\
Ga$_{3}$Hg & mp-1184289 & P6\_3/mmc & 0.2216 & -1.4162 & 0.0000 & 0.0000 \\
Ga$_{3}$Pb & mp-1184425 & I4/mmm & 0.1929 & -1.0646 & 0.0042 & 0.0000 \\
Ga$_{3}$Re\textcolor{red}{$^*$} & mp-1188393 & P4\_2/mnm & 0.4857 & -2.3451 & 0.0000 & 0.0000 \\
GaSb & mp-1224785 & I-4m2 & 0.2847 & -1.4795 & 0.0030 & 0.0000 \\
Ge$^*$ & mp-32 & Fd-3m & 0.8956 & -4.6450 & 0.0015 & 0.0000 \\
Ge$_{2}$Os$^*$ & mp-10032 & C2/m & 0.4042 & -2.1066 & 0.0069 & 0.0003 \\
Ge$_{3}$Mo$_{5}$\textcolor{red}{$^*$} & mp-17094 & I4/mcm & 0.4229 & -3.3674 & 0.0068 & 0.0009 \\
Ge$_{3}$O & mp-1184665 & I4/mmm & 0.1453 & -0.8722 & 0.0091 & 0.0027 \\
Ge$_{4}$Rh\textcolor{red}{$^*$} & mp-1104286 & P3\_121 & 0.3355 & -1.8036 & 0.0034 & 0.0006 \\
GeAs\textcolor{red}{$^*$} & mp-7591 & I4mm & 0.4634 & -2.4402 & 0.0063 & 0.0014 \\
GeSb & mp-9935 & Fm-3m & 0.4373 & -2.1513 & 0.0101 & 0.0039 \\
H$_{3}$Pd$_{4}$ & mp-1224413 & R-3m & 0.1466 & -1.0855 & 0.1005 & 0.0202 \\
H$_{4}$N & mp-1238814 & P-43m & 0.6754 & -6.0262 & 0.0498 & 0.0006 \\
HC$_{2}$ & mp-995184 & C2/m & 1.2737 & -11.4891 & 0.0000 & 0.0025 \\
HC$_{3}$ & mp-995200 & P2\_1/m & 1.2493 & -11.5739 & 0.0000 & 0.0000 \\
HC$_{6}$ & mp-995202 & P-1 & 1.2320 & -11.5624 & 0.0000 & 0.0000 \\
Hf$^*$ & mp-103 & P6\_3/mmc & 0.3121 & -1.1612 & 0.0000 & 0.0000 \\
Hf$_{3}$Pb & mp-1184772 & I4/mmm & 0.3286 & -1.3613 & 0.0000 & 0.0000 \\
HfCd$^*$ & mp-1007758 & P4/nmm & 0.3741 & -1.6812 & 0.0000 & 0.0000 \\
HfCr$_{2}$\textcolor{red}{$^*$} & mp-569624 & P6\_3/mmc & 0.4565 & -1.1713 & 0.0000 & 0.0000 \\
HfCu$_{3}$ & mp-1184704 & I4/mmm & 0.1762 & -0.7182 & 0.0196 & 0.0000 \\
HfMg$_{15}$ & mp-1023478 & P-6m2 & 0.0440 & -0.1164 & 1.0039 & 0.4374 \\
HfMo & mp-1224314 & Cmmm & 0.3368 & -1.2722 & 0.0000 & 0.0000 \\
HfTc$_{2}$$^*$ & mp-1095669 & P6\_3/mmc & 0.3775 & -1.4429 & 0.0098 & 0.0012 \\
HfW$_{2}$ & mp-1224230 & Fmmm & 0.4943 & -2.5963 & 0.0000 & 0.0000 \\
Hg\textcolor{red}{$^*$} & mp-10861 & P6/mmm & 0.1622 & -0.9652 & 0.0000 & 0.0000 \\
Hg$_{3}$Pb & mp-1184700 & P6\_3/mmc & 0.2256 & -1.4191 & 0.0000 & 0.0000 \\
Hg$_{3}$Sb & mp-977242 & P6\_3/mmc & 0.2146 & -1.3610 & 0.0001 & 0.0003 \\
Hg$_{4}$Pt$^*$ & mp-936 & Im-3m & 0.5477 & -3.5588 & 0.0000 & 0.0005 \\
Hg$_{5}$Pd$_{2}$\textcolor{red}{$^*$} & mp-30721 & P4/mbm & 0.2364 & -1.7271 & 0.0000 & 0.0002 \\
IN$_{2}$ & mp-1097011 & I4/mcm & 2.9506 & -23.1289 & 0.0177 & 0.0275 \\
In\textcolor{red}{$^*$} & mp-1055994 & I4/mmm & 0.2177 & -1.1562 & 0.0000 & 0.0000 \\
In$_{2}$Te$_{3}$\textcolor{red}{$^*$} & mp-1105025 & I4mm & 0.6032 & -2.6916 & 0.0224 & 0.0090 \\
In$_{3}$Ir\textcolor{red}{$^*$} & mp-630976 & P4\_2/mnm & 0.5910 & -3.4701 & 0.0000 & 0.0000 \\
In$_{7}$Cu$_{10}$$^*$ & mp-646039 & C2/m & 0.1416 & -0.8072 & 0.0230 & 0.0039 \\
InNi\textcolor{red}{$^*$} & mp-20997 & Pm-3m & 0.2785 & -1.7672 & 0.0081 & 0.0156 \\
InSe$_{2}$ & mp-1232036 & P4/nmm & 0.2771 & -1.4780 & 0.0349 & 0.0038 \\
InSn$_{4}$ & mp-1223653 & Cmmm & 0.3867 & -2.1185 & 0.0000 & 0.0000 \\
Ir$^*$ & mp-101 & Fm-3m & 0.3469 & -2.4130 & 0.0317 & 0.0274 \\
Ir$_{3}$Ru & mp-974358 & I4/mmm & 0.3767 & -2.6659 & 0.0317 & 0.0198 \\
Ir$_{4}$W & mp-1223666 & R-3m & 0.3980 & -2.8035 & 0.0422 & 0.0434 \\
IrPd & mp-1223675 & R-3m & 0.3508 & -2.9933 & 0.0756 & 0.0384 \\
K\textcolor{red}{$^*$} & mp-10157 & Fm-3m & 0.0430 & -0.0808 & 0.0923 & 0.0499 \\
K$_{2}$O & mp-1223784 & P4/mmm & 0.1375 & -0.3763 & 0.0831 & 0.0000 \\
K$_{3}$Rb & mp-976115 & P6\_3/mmc & 0.0431 & -0.0882 & 0.1130 & 0.0889 \\
K$_{5}$As$_{4}$ & mp-684799 & P-1 & 0.9859 & -3.9306 & 0.0452 & 0.0189 \\
KMg$_{15}$ & mp-1023498 & P-6m2 & 0.0006 & -0.0094 & 4.8845 & 1.9656 \\
KN$_{2}$\textcolor{red}{$^*$} & mp-1071868 & Cmc2\_1 & 1.0490 & -7.9889 & 0.0202 & 0.0000 \\
KN$_{2}$ & mp-1180757 & Cmcm & 1.0210 & -7.4152 & 0.0411 & 0.0000 \\
KN$_{3}$ & mp-636056 & I4/mcm & 2.8711 & -22.8747 & 0.0090 & 0.0107 \\
KP\textcolor{red}{$^*$} & mp-1058315 & R-3m & 0.1254 & -0.4533 & 0.0190 & 0.0442 \\
KP & mp-1057787 & Immm & 0.1070 & -0.4621 & 0.0170 & 0.0218 \\
KSn$_{3}$ & mp-1185151 & P6\_3/mmc & 0.4006 & -1.5431 & 0.0179 & 0.0027 \\
Li$^*$ & mp-1018134 & R-3m & 0.0886 & -0.6931 & 0.0000 & 0.0000 \\
Li$_{13}$Sn$_{5}$\textcolor{red}{$^*$} & mp-30769 & P-3m1 & 0.6069 & -2.7881 & 0.0000 & 0.0000 \\
Li$_{3}$Ag & mp-976408 & I4/mmm & 0.1084 & -1.0444 & 0.0000 & 0.0000 \\
Li$_{3}$Ca & mp-976075 & I4/mmm & 0.1652 & -0.6223 & 0.0000 & 0.0000 \\
Li$_{3}$Hg\textcolor{red}{$^*$} & mp-1646 & Fm-3m & 0.1488 & -1.2194 & 0.0000 & 0.0000 \\
LiB$^*$ & mp-1001835 & P6\_3/mmc & 1.1571 & -9.9815 & 0.0000 & 0.0000 \\
LiRh$_{3}$$^*$ & mp-1001787 & Imm2 & 0.1776 & -1.5461 & 0.0814 & 0.0192 \\
LiTl$_{3}$ & mp-1185253 & I4/mmm & 0.2370 & -1.2989 & 0.0000 & 0.0000 \\
Mg$^*$ & mp-153 & P6\_3/mmc & 0.0000 & 0.0000 & -- & -- \\
Mg$_{15}$B & mp-1023515 & P-6m2 & 0.0280 & -0.1152 & 0.9423 & 0.5909 \\
Mg$_{15}$Ni & mp-1023622 & P-6m2 & 0.0020 & -0.0213 & 3.2054 & 2.1119 \\
Mg$_{2}$Ge\textcolor{red}{$^*$} & mp-250 & Fd-3m & 0.1806 & -0.5039 & 0.1144 & 0.0306 \\
Mg$_{2}$Ni$^*$ & mp-2137 & P6\_222 & 0.0818 & -0.2478 & 0.6457 & 0.3329 \\
Mg$_{3}$Al & mp-978271 & Pm-3m & 0.0812 & -0.3224 & 0.1540 & 0.0681 \\
Mg$_{3}$Ru$_{2}$$^*$ & mp-625 & P4\_132 & 0.5409 & -1.6132 & 0.1816 & 0.0936 \\
Mg$_{3}$Zr & mp-1094466 & P6\_3/mmc & 0.2889 & -0.3975 & 0.0887 & 0.0141 \\
Mg$_{5}$Sc & mp-1185837 & R32 & 0.1000 & -0.1246 & 0.2961 & 0.0068 \\
Mg$_{5}$Te & mp-1185828 & R32 & 0.0108 & -0.0917 & 1.0311 & 0.5406 \\
MgAl$_{2}$\textcolor{red}{$^*$} & mp-1102064 & I4\_1/amd & 0.3363 & -1.6067 & 0.0037 & 0.0000 \\
MgCd$^*$ & mp-11307 & Pmma & 0.4362 & -2.0470 & 0.0026 & 0.0006 \\
MgIn$_{2}$\textcolor{red}{$^*$} & mp-1077016 & Fd-3m & 0.4220 & -1.8995 & 0.0118 & 0.0002 \\
MgPb$_{2}$ & mp-2016042 & C2/m & 0.2720 & -1.0252 & 0.0628 & 0.0086 \\
MgPd$_{3}$\textcolor{red}{$^*$} & mp-7753 & Pm-3m & 0.2542 & -1.7837 & 0.0264 & 0.0304 \\
MgSn & mp-1094669 & Pm-3m & 0.2238 & -0.7050 & 0.1000 & 0.0334 \\
MgSn & mp-1094801 & P4/mmm & 0.3134 & -1.2760 & 0.1099 & 0.0281 \\
MgTi & mp-1094346 & P2\_1/m & 0.4231 & -0.7659 & 0.0009 & 0.0000 \\
MgTl$_{3}$ & mp-978262 & P6\_3/mmc & 0.2760 & -1.3397 & 0.0002 & 0.0000 \\
Mn$^*$ & mp-35 & I-43m & -- & -0.6082 & 0.0411 & 0.0047 \\
Mn$_{4}$B & mp-1221769 & F222 & 0.3919 & -2.3186 & 0.0085 & 0.0007 \\
Mn$_{4}$Ga & mp-1221812 & R-3m & 0.3290 & -2.4448 & 0.0073 & 0.0017 \\
MnAl$_{12}$$^*$ & mp-1104165 & Im-3 & 0.2798 & -1.6283 & 0.0000 & 0.0000 \\
MnO$_{2}$\textcolor{red}{$^*$} & mp-510408 & P4\_2/mnm & 0.4433 & -3.0527 & 0.0548 & 0.0603 \\
Mo$^*$ & mp-129 & Im-3m & 0.3113 & -1.1059 & 0.0139 & 0.0005 \\
Mo$_{2}$C & mp-1221498 & P-3m1 & 0.5408 & -3.5050 & 0.0057 & 0.0002 \\
Mo$_{3}$Se$_{4}$\textcolor{red}{$^*$} & mp-21021 & R-3 & 0.4924 & -2.4124 & 0.0151 & 0.0004 \\
MoPd$_{2}$ & mp-1206704 & Immm & 0.5227 & -1.7521 & 0.0013 & 0.0000 \\
MoRu & mp-1221386 & P-6m2 & 0.3772 & -1.3952 & 0.0182 & 0.0004 \\
N$_{2}$ & mp-1059834 & Imma & 1.3533 & -15.9343 & 0.0159 & 0.0010 \\
NO & mp-1180334 & P2\_12\_12\_1 & 1.8010 & -21.3123 & 0.0427 & 0.0039 \\
Na$^*$ & mp-10172 & P6\_3/mmc & 0.0470 & -0.0839 & 0.1168 & 0.0901 \\
Na$_{2}$Cl & mp-1077015 & Cmmm & 0.0377 & -0.1951 & 0.0901 & 0.0029 \\
Na$_{2}$Cl & mp-990084 & R-3m & 0.0479 & -0.3527 & 0.0618 & 0.0000 \\
Na$_{3}$Cl & mp-1064484 & P4/mmm & 0.0297 & -0.2254 & 0.0859 & 0.0249 \\
Na$_{3}$Hg & mp-1186050 & I4/mmm & 0.1243 & -0.3788 & 0.0010 & 0.0000 \\
Na$_{3}$In & mp-864758 & I4/mmm & 0.1784 & -0.3865 & 0.0055 & 0.0000 \\
Na$_{3}$Mg & mp-982637 & I4/mmm & 0.0436 & -0.0860 & 0.1585 & 0.0598 \\
Na$_{3}$Si$_{34}$ & mp-1221235 & R-3m & 0.8887 & -4.3194 & 0.0033 & 0.0052 \\
Nb$^*$ & mp-75 & Im-3m & 0.3268 & -1.1439 & 0.0000 & 0.0000 \\
Nb$_{3}$Ir$^*$ & mp-1458 & Pm-3n & 0.4144 & -2.0103 & 0.0000 & 0.0002 \\
Nb$_{3}$Ni & mp-999394 & I4/mmm & 0.4256 & -1.5440 & 0.0000 & 0.0002 \\
Nb$_{4}$Rh & mp-1220441 & Fmmm & 0.3405 & -1.1843 & 0.0002 & 0.0000 \\
NbCd$_{3}$\textcolor{red}{$^*$} & mp-11309 & Pm-3m & 0.2671 & -1.6744 & 0.0000 & 0.0000 \\
NbGa$_{3}$$^*$ & mp-1973 & I4/mmm & 0.3104 & -1.7262 & 0.0000 & 0.0000 \\
NbO\textcolor{red}{$^*$} & mp-2692 & Fm-3m & 0.2933 & -2.2812 & 0.0292 & 0.0023 \\
NbO & mp-634965 & P4/mmm & 0.9184 & -5.5839 & 0.0000 & 0.0000 \\
NbP\textcolor{red}{$^*$} & mp-9830 & I4\_1/amd & 0.3299 & -1.7803 & 0.0001 & 0.0000 \\
NbP & mp-1220317 & Cmmm & 0.3238 & -1.7805 & 0.0084 & 0.0004 \\
NbRh$_{2}$ & mp-1220414 & Pm & 0.3585 & -1.4459 & 0.0440 & 0.0023 \\
NbS$_{2}$ & mp-774873 & P6\_3/mmc & 0.6390 & -3.1946 & 0.0049 & 0.0000 \\
NbZn$_{16}$$^*$ & mp-640041 & Cmcm & 0.2534 & -1.6346 & 0.0000 & 0.0000 \\
Ni$^*$ & mp-23 & Fm-3m & 0.0947 & -0.4204 & 0.3799 & 0.2037 \\
Ni$_{3}$Mo$^*$ & mp-11506 & Pmmn & 0.2915 & -1.0002 & 0.1370 & 0.0141 \\
NiB\textcolor{red}{$^*$} & mp-14019 & Cmcm & 0.6996 & -6.1857 & 0.0198 & 0.0093 \\
NiHg$_{4}$\textcolor{red}{$^*$} & mp-570835 & Im-3m & 0.3125 & -2.2179 & 0.0000 & 0.0004 \\
NiS\textcolor{red}{$^*$} & mp-1547 & R3m & 0.2630 & -1.5895 & 0.1140 & 0.0342 \\
NiSb$_{3}$\textcolor{red}{$^*$} & mp-1106023 & Im-3 & 0.3169 & -1.4693 & 0.0547 & 0.0298 \\
Os$^*$ & mp-49 & P6\_3/mmc & 0.2567 & -1.1579 & 0.0538 & 0.0135 \\
OsC\textcolor{red}{$^*$} & mp-7142 & P-6m2 & 0.4688 & -3.0941 & 0.0402 & 0.0068 \\
OsC & mp-1009537 & Pm-3m & 0.3353 & -2.0943 & 0.0446 & 0.0000 \\
OsC & mp-1009539 & Fm-3m & 0.4618 & -3.1760 & 0.0828 & 0.0058 \\
OsC & mp-999308 & P6\_3/mmc & 0.4521 & -3.1583 & 0.0525 & 0.0226 \\
PW$_{3}$ & mp-1209828 & I-42m & 0.4114 & -2.7903 & 0.0000 & 0.0000 \\
Pb$^*$ & mp-20483 & Fm-3m & 0.1943 & -0.8519 & 0.0276 & 0.0000 \\
Pb$_{2}$Au$^*$ & mp-22795 & I4/mcm & 0.4524 & -3.2812 & 0.0000 & 0.0000 \\
PbAu$_{2}$\textcolor{red}{$^*$} & mp-19871 & Fd-3m & 0.3048 & -2.3586 & 0.0000 & 0.0000 \\
PbAu$_{3}$ & mp-982246 & Fm-3m & 0.1941 & -1.4564 & 0.0112 & 0.0022 \\
Pd$^*$ & mp-2 & Fm-3m & 0.1631 & -1.2162 & 0.1264 & 0.0877 \\
Pt$^*$ & mp-126 & Fm-3m & 0.2367 & -1.8089 & 0.0927 & 0.0819 \\
Rb\textcolor{red}{$^*$} & mp-70 & Im-3m & 0.0590 & -0.1180 & 0.1307 & 0.0377 \\
Rb$_{2}$Hg$_{7}$\textcolor{red}{$^*$} & mp-31474 & P-3m1 & 0.1920 & -0.9499 & 0.0000 & 0.0000 \\
Rb$_{3}$In & mp-974948 & Pm-3m & 0.1129 & -0.2266 & 0.0341 & 0.0000 \\
Rb$_{3}$Li & mp-974938 & Pm-3m & 0.0991 & -0.3104 & 0.0185 & 0.0000 \\
Rb$_{3}$P$_{11}$ & mp-1173444 & C2/m & 0.8390 & -4.2144 & 0.0355 & 0.0027 \\
RbP\textcolor{red}{$^*$} & mp-1179688 & C2/m & 0.1031 & -0.4288 & 0.0138 & 0.0467 \\
RbP & mp-1058499 & Immm & 0.1277 & -0.4304 & 0.0099 & 0.0324 \\
Re\textcolor{red}{$^*$} & mp-8 & P6\_3/mmc & 0.3162 & -1.3803 & 0.0205 & 0.0030 \\
Re$_{2}$Si\textcolor{red}{$^*$} & mp-669547 & P2\_1/c & 0.3849 & -1.8091 & 0.0190 & 0.0006 \\
Re$_{3}$Ge$_{7}$\textcolor{red}{$^*$} & mp-29141 & Cmcm & 0.4553 & -2.2942 & 0.0022 & 0.0002 \\
Re$_{3}$N$_{5}$ & mp-1080194 & Imm2 & 0.8324 & -5.2882 & 0.0269 & 0.0001 \\
Re$_{3}$Pd & mp-862586 & P6\_3/mmc & 0.3788 & -1.7449 & 0.0020 & 0.0011 \\
Re$_{5}$Si$_{3}$\textcolor{red}{$^*$} & mp-1105959 & I4/mcm & 0.3379 & -2.7922 & 0.0251 & 0.0018 \\
ReC\textcolor{red}{$^*$} & mp-1079494 & P6\_3/mmc & 0.4741 & -3.3056 & 0.0211 & 0.0022 \\
ReC & mp-1009731 & Pm-3m & 0.3588 & -2.2826 & 0.0149 & 0.0003 \\
Rh$^*$ & mp-74 & Fm-3m & 0.1620 & -0.7081 & 0.1852 & 0.0924 \\
Ru$^*$ & mp-33 & P6\_3/mmc & 0.2375 & -0.9478 & 0.0875 & 0.0211 \\
RuC\textcolor{red}{$^*$} & mp-10030 & P-6m2 & 0.4770 & -2.9147 & 0.0418 & 0.0078 \\
RuC & mp-1009787 & Fm-3m & 0.4752 & -2.9734 & 0.0862 & 0.0067 \\
RuRh$_{3}$ & mp-974335 & Pm-3m & 0.2074 & -0.8611 & 0.1382 & 0.0326 \\
SN & mp-1078816 & P2\_1/c & 0.2823 & -1.7903 & 0.0359 & 0.0006 \\
SN & mp-1173453 & P2\_1 & 1.5195 & -10.9244 & 0.0458 & 0.0000 \\
SN & mp-684642 & P2\_1 & 1.3504 & -13.4315 & 0.0333 & 0.0038 \\
Sb$^*$ & mp-104 & R-3m & 0.4398 & -1.7202 & 0.0425 & 0.0063 \\
Sb$_{2}$S$_{5}$ & mp-1359089 & P2\_1/c & 1.0848 & -6.8891 & 0.0690 & 0.0097 \\
Sc$^*$ & mp-67 & P6\_3/mmc & 0.2110 & -0.6596 & 0.0000 & 0.0000 \\
Sc$_{5}$Ge$_{3}$$^*$ & mp-17190 & P6\_3/mcm & 0.3162 & -1.5145 & 0.0000 & 0.0000 \\
ScAg$_{2}$$^*$ & mp-1018128 & I4/mmm & 0.2014 & -1.5516 & 0.0000 & 0.0000 \\
ScOs$_{3}$ & mp-981384 & Pm-3m & 0.2618 & -1.1903 & 0.0252 & 0.0106 \\
ScTc$_{3}$ & mp-867262 & Fm-3m & 0.3043 & -1.0739 & 0.0080 & 0.0010 \\
ScTe$^*$ & mp-10026 & P6\_3/mmc & 0.3365 & -1.5673 & 0.0000 & 0.0000 \\
Si$_{3}$Pt & mp-972763 & P6\_3/mmc & 0.4284 & -2.6580 & 0.0034 & 0.0000 \\
Si$_{3}$W$_{5}$\textcolor{red}{$^*$} & mp-31219 & I4/mcm & 0.4312 & -2.5598 & 0.0000 & 0.0000 \\
SiAu$_{3}$ & mp-1186998 & Fm-3m & 0.1974 & -1.4121 & 0.0156 & 0.0000 \\
SiO$_{2}$ & mp-556044 & P-31c & 0.7437 & -6.4399 & 0.0117 & 0.0000 \\
SiO$_{2}$ & mp-638049 & F4\_132 & 0.6594 & -6.0556 & 0.0488 & 0.0000 \\
SiRh\textcolor{red}{$^*$} & mp-2287 & P2\_1/c & 0.2801 & -1.6004 & 0.0243 & 0.0009 \\
SiTc$_{3}$ & mp-972784 & P6\_3/mmc & 0.3385 & -1.4573 & 0.0257 & 0.0020 \\
SiTe$_{2}$\textcolor{red}{$^*$} & mp-2755 & P-3m1 & 0.5658 & -2.7523 & 0.0695 & 0.0130 \\
Sn\textcolor{red}{$^*$} & mp-623511 & I4/mmm & 0.1989 & -0.8803 & 0.0235 & 0.0006 \\
Sn$^*$ & mp-117 & Fd-3m & 0.8707 & -3.6930 & 0.0001 & 0.0004 \\
Sn$_{3}$S & mp-1187038 & Fm-3m & 0.2241 & -1.1171 & 0.0176 & 0.0000 \\
Sn$_{3}$Se & mp-1187055 & Fm-3m & 0.2108 & -1.0119 & 0.0214 & 0.0027 \\
Sn$_{7}$Ir$_{5}$\textcolor{red}{$^*$} & mp-20466 & I4mm & 0.5375 & -3.2800 & 0.0000 & 0.0000 \\
Sn$_{7}$Pd$_{5}$ & mp-1219006 & C2 & 0.3765 & -2.7926 & 0.0000 & 0.0004 \\
SnP\textcolor{red}{$^*$} & mp-475 & Fm-3m & 0.4347 & -2.3226 & 0.0105 & 0.0058 \\
SnPb$_{3}$ & mp-1187118 & P6\_3/mmc & 0.1988 & -0.8714 & 0.0218 & 0.0000 \\
SnPt$^*$ & mp-19856 & P6\_3/mmc & 0.4612 & -3.2535 & 0.0000 & 0.0000 \\
Sr$^*$ & mp-139 & P6\_3/mmc & 0.0009 & -0.0171 & 4.9050 & 5.3659 \\
Sr$_{2}$As\textcolor{red}{$^*$} & mp-15697 & I4/mmm & 0.0787 & -0.2746 & 0.2097 & 0.0267 \\
Sr$_{3}$N$_{2}$ & mp-1013529 & Pm-3m & 0.1064 & -0.4148 & 0.0858 & 0.0000 \\
Sr$_{3}$Sn & mp-1187136 & Pm-3m & 0.0501 & -0.2374 & 0.4377 & 0.0368 \\
Sr$_{5}$Sn$_{3}$\textcolor{red}{$^*$} & mp-17720 & I4/mcm & 0.8182 & -2.1022 & 0.1796 & 0.0000 \\
SrAl$_{2}$\textcolor{red}{$^*$} & mp-1638 & Fd-3m & 0.4263 & -1.6649 & 0.0004 & 0.0000 \\
SrB$_{2}$ & mp-1008755 & P6/mmm & 0.9937 & -5.2175 & 0.0052 & 0.0007 \\
SrN\textcolor{red}{$^*$} & mp-1058586 & Fm-3m & 0.1230 & -0.6070 & 0.0462 & 0.0849 \\
SrO$_{10}$ & mp-1208655 & P4/mmm & 0.8181 & -6.3571 & 0.1157 & 0.0000 \\
SrSi$^*$ & mp-2661 & Cmcm & 0.8077 & -2.8883 & 0.0564 & 0.0097 \\
SrTl$_{3}$ & mp-1206636 & Pm-3m & 0.2857 & -1.2068 & 0.0201 & 0.0007 \\
Ta\textcolor{red}{$^*$} & mp-50 & Im-3m & 0.3226 & -1.2794 & 0.0000 & 0.0000 \\
Ta$_{4}$C$_{3}$ & mp-1218000 & R3m & 0.5265 & -3.5440 & 0.0000 & 0.0000 \\
Ta$_{5}$Ge$_{3}$$^*$ & mp-17593 & I4/mcm & 0.3556 & -2.6739 & 0.0000 & 0.0000 \\
Ta$_{7}$Co$_{6}$\textcolor{red}{$^*$} & mp-1104548 & R-3m & 0.4891 & -3.0211 & 0.0018 & 0.0000 \\
TaBe$_{12}$$^*$ & mp-567842 & I4/mmm & 0.2080 & -1.8875 & 0.0000 & 0.0000 \\
TaN\textcolor{red}{$^*$} & mp-1497 & P6/mmm & 0.4241 & -2.9806 & 0.0009 & 0.0000 \\
TaN & mp-1009831 & Pm-3m & 0.3185 & -2.1377 & 0.0002 & 0.0002 \\
TaN & mp-1009833 & F-43m & 0.6451 & -4.4476 & 0.0000 & 0.0006 \\
TaP$^*$ & mp-1067587 & I4\_1md & 0.4762 & -2.7480 & 0.0000 & 0.0000 \\
TaP & mp-1187244 & P-6m2 & 0.4792 & -2.7509 & 0.0000 & 0.0000 \\
TaRe & mp-1217894 & Cmmm & 0.3311 & -1.5114 & 0.0026 & 0.0003 \\
TaSb$_{2}$$^*$ & mp-11697 & C2/m & 0.4899 & -2.0334 & 0.0002 & 0.0000 \\
Tc$^*$ & mp-113 & P6\_3/mmc & 0.3130 & -1.2061 & 0.0240 & 0.0006 \\
Tc$_{3}$Mo & mp-1187395 & Pm-3m & 0.3332 & -1.2727 & 0.0068 & 0.0014 \\
Tc$_{3}$Os & mp-862779 & P6\_3/mmc & 0.3182 & -1.2186 & 0.0252 & 0.0052 \\
Tc$_{3}$W & mp-1187398 & P6\_3/mmc & 0.4723 & -2.3705 & 0.0000 & 0.0000 \\
TcW & mp-1217337 & Cmmm & 0.6035 & -3.6836 & 0.0000 & 0.0002 \\
Te$_{2}$Au\textcolor{red}{$^*$} & mp-567525 & C2/m & 0.5036 & -3.0627 & 0.0497 & 0.0327 \\
Te$_{2}$Ir\textcolor{red}{$^*$} & mp-1551 & C2/m & 0.7621 & -4.4689 & 0.0230 & 0.0046 \\
Te$_{2}$Pd$^*$ & mp-782 & P-3m1 & 0.5813 & -3.5060 & 0.0291 & 0.0200 \\
Te$_{2}$Pd$_{7}$ & mp-1217443 & R-3m & 0.2769 & -1.8731 & 0.0363 & 0.0145 \\
Ti\textcolor{red}{$^*$} & mp-46 & P6\_3/mmc & 0.3010 & -0.7536 & 0.0000 & 0.0000 \\
Ti$_{2}$Cu$^*$ & mp-742 & I4/mmm & 0.3374 & -0.7591 & 0.0016 & 0.0000 \\
Ti$_{3}$Ge\textcolor{red}{$^*$} & mp-1189044 & I-4 & 0.3243 & -1.6051 & 0.0000 & 0.0000 \\
Ti$_{3}$In$^*$ & mp-866046 & P6\_3/mmc & 0.3128 & -1.1764 & 0.0000 & 0.0000 \\
Ti$_{3}$Nb & mp-1217091 & Amm2 & 0.3774 & -1.2750 & 0.0000 & 0.0000 \\
Ti$_{3}$Pt$^*$ & mp-2339 & Pm-3n & 0.3632 & -1.7698 & 0.0000 & 0.0000 \\
TiBe$_{12}$$^*$ & mp-1104067 & I4/mmm & 0.2186 & -1.9904 & 0.0000 & 0.0000 \\
TiMo$_{2}$ & mp-1216675 & Fmmm & 0.3395 & -1.1740 & 0.0000 & 0.0000 \\
TiMo$_{2}$ & mp-1216675 & Fmmm & 0.3395 & -1.1740 & 0.0000 & 0.0000 \\
TiNi$^*$ & mp-1048 & P2\_1/m & 0.2818 & -0.8685 & 0.0059 & 0.0000 \\
TiNi & mp-1067248 & Cmcm & 0.2798 & -0.8586 & 0.0060 & 0.0000 \\
TiO$^*$ & mp-1071163 & P-62m & 0.3594 & -2.4547 & 0.0015 & 0.0000 \\
TiSi$_{2}$$^*$ & mp-1077503 & Cmcm & 0.4599 & -3.0700 & 0.0000 & 0.0000 \\
TiSi$_{2}$ & mp-570991 & Immm & 0.2659 & -1.2939 & 0.0000 & 0.0000 \\
TiW & mp-1216621 & Cmmm & 0.6568 & -3.6207 & 0.0000 & 0.0000 \\
Tl$^*$ & mp-82 & P6\_3/mmc & 0.1905 & -1.0463 & 0.0000 & 0.0000 \\
Tl$_{3}$Au & mp-1080569 & Pm-3n & 0.2291 & -1.2850 & 0.0000 & 0.0000 \\
Tl$_{3}$Ga & mp-1187539 & Fm-3m & 0.2090 & -1.0914 & 0.0000 & 0.0000 \\
Tl$_{4}$Sn & mp-1216560 & R-3m & 0.2039 & -1.3837 & 0.0010 & 0.0008 \\
TlAs & mp-1007770 & F-43m & 0.7362 & -3.5581 & 0.0000 & 0.0002 \\
TlAu$_{3}$ & mp-1187508 & Pm-3m & 0.2130 & -1.5377 & 0.0000 & 0.0000 \\
TlC\textcolor{red}{$^*$} & mp-567465 & Fm-3m & 0.3776 & -2.4265 & 0.0000 & 0.0010 \\
TlC & mp-1058887 & P4/mmm & 0.4837 & -3.1370 & 0.0000 & 0.0003 \\
TlN\textcolor{red}{$^*$} & mp-1017982 & P6\_3mc & 0.6066 & -3.6824 & 0.0023 & 0.0000 \\
TlN & mp-1002222 & Fm-3m & 0.3446 & -2.1516 & 0.0029 & 0.0000 \\
TlN & mp-1007778 & F-43m & 0.6105 & -3.6815 & 0.0017 & 0.0000 \\
V$^*$ & mp-146 & Im-3m & 0.3178 & -0.8777 & 0.0000 & 0.0000 \\
VCo$_{3}$ & mp-1095057 & P6\_3/mmc & 0.3730 & -2.5141 & 0.0012 & 0.0023 \\
VN & mp-1017532 & P6\_3/mmc & 0.4293 & -2.6954 & 0.0000 & 0.0000 \\
VNi$_{2}$$^*$ & mp-11531 & Immm & 0.4769 & -1.4029 & 0.0918 & 0.0008 \\
W$^*$ & mp-91 & Im-3m & 0.4372 & -2.5600 & 0.0000 & 0.0001 \\
W$_{9}$C$_{4}$ & mp-684989 & R32 & 0.6031 & -4.5942 & 0.0000 & 0.0000 \\
WBr$_{5}$ & mp-1436789 & P-1 & 0.8263 & -4.3695 & 0.0774 & 0.0061 \\
WC\textcolor{red}{$^*$} & mp-13136 & Fm-3m & 0.5508 & -4.2240 & 0.0012 & 0.0009 \\
WC & mp-1008630 & Pm-3m & 0.4127 & -2.9246 & 0.0000 & 0.0010 \\
WC & mp-1008635 & F-43m & 0.8710 & -6.2790 & 0.0000 & 0.0000 \\
WO$_{2}$$^*$ & mp-1524394 & P2\_1/c & 0.6713 & -5.7408 & 0.0039 & 0.0000 \\
Y$^*$ & mp-112 & P6\_3/mmc & 0.2374 & -0.8682 & 0.0000 & 0.0000 \\
Y$_{3}$Co$^*$ & mp-1105598 & Pnma & 0.4720 & -2.6094 & 0.0000 & 0.0000 \\
Y$_{3}$Co & mp-1189329 & P2\_1 & 0.4718 & -2.6045 & 0.0000 & 0.0000 \\
YAg$_{3}$ & mp-1187811 & Fm-3m & 0.1780 & -1.2749 & 0.0000 & 0.0000 \\
Zn$^*$ & mp-79 & P6\_3/mmc & 0.2296 & -1.8892 & 0.0000 & 0.0000 \\
Zn & mp-2647117 & Im-3m & 0.1819 & -1.3641 & 0.0000 & 0.0000 \\
Zn$_{13}$Rh$^*$ & mp-13447 & C2/m & 0.1989 & -1.5615 & 0.0000 & 0.0000 \\
Zn$_{2}$Ag & mp-1093999 & P6\_3/mmc & 0.2898 & -2.4762 & 0.0000 & 0.0000 \\
Zn$_{3}$Au & mp-1188000 & P6\_3/mmc & 0.2275 & -1.8469 & 0.0000 & 0.0000 \\
Zn$_{3}$N$_{2}$ & mp-1524551 & P6\_3/mmc & 0.3888 & -2.8616 & 0.0030 & 0.0000 \\
Zn$_{3}$N$_{2}$$^*$ & mp-9460 & Ia-3 & 0.4180 & -3.1067 & 0.0036 & 0.0000 \\
Zn$_{7}$Tc & mp-1206378 & Fm-3m & 0.2749 & -1.7542 & 0.0000 & 0.0009 \\
ZnAg\textcolor{red}{$^*$} & mp-1912 & Pm-3m & 0.2035 & -1.8265 & 0.0000 & 0.0000 \\
ZnPb$_{3}$ & mp-1187949 & Fm-3m & 0.2392 & -1.3227 & 0.0004 & 0.0004 \\
ZnSn & -- & P6\_3/mmc & 0.4604 & -3.0247 & 0.0000 & 0.0000 \\
Zr$^*$ & mp-131 & P6\_3/mmc & 0.3223 & -1.3170 & 0.0000 & 0.0000 \\
Zr$_{2}$Au$^*$ & mp-2348 & I4/mmm & 0.3095 & -2.0203 & 0.0000 & 0.0000 \\
Zr$_{3}$Be & mp-1188016 & I4/mmm & 0.3655 & -1.7390 & 0.0000 & 0.0000 \\
Zr$_{3}$Fe$^*$ & mp-31205 & Cmcm & 0.4128 & -1.4259 & 0.0000 & 0.0000 \\
Zr$_{3}$Ge & mp-1188039 & P6\_3/mmc & 0.3460 & -1.7196 & 0.0000 & 0.0000 \\
Zr$_{5}$Pb$_{3}$$^*$ & mp-681992 & P6\_3/mcm & 0.3301 & -1.9739 & 0.0000 & 0.0000 \\
ZrBr\textcolor{red}{$^*$} & mp-1065846 & C2/m & 0.3528 & -2.0185 & 0.0000 & 0.0001 \\
ZrGa$_{3}$\textcolor{red}{$^*$} & mp-30685 & I4/mmm & 0.2933 & -1.8532 & 0.0000 & 0.0000 \\
ZrHg$_{3}$\textcolor{red}{$^*$} & mp-1608 & Pm-3m & 0.2490 & -1.9086 & 0.0000 & 0.0000 \\
ZrMo$_{2}$$^*$ & mp-2049 & Fd-3m & 0.4671 & -1.6863 & 0.0002 & 0.0000 \\
ZrRh\textcolor{red}{$^*$} & mp-2808 & Pm-3m & 0.3062 & -1.4732 & 0.0094 & 0.0024 \\
ZrTi & mp-1215236 & P4/mmm & 0.3777 & -1.6592 & 0.0000 & 0.0000 \\
\end{longtable}

}

\subsection{Ionic Compounds}
\label{sec:si-master-ionic}
{\footnotesize
\input{si_master_list_ionic_en.tex}
}

\subsection{Covalent Compounds}
\label{sec:si-master-covalent}
{\footnotesize
\input{si_master_list_covalent_en.tex}
}

\subsection{Mixed (Zintl-Type) Compounds}
\label{sec:si-master-mixed}
{\footnotesize
\begin{longtable}{@{}lllrrrr@{}}
\caption{Mixed-bonding (Zintl-type) compounds (a strongly covalent homoatomic pair coexists with a markedly weaker heteroatomic pair, e.g. the azide N$_3^-$ ion ionically bound to its cation; observed range 0.8490 to 2.7322) --- same columns.}\\
\toprule
Formula & mp-id & Space group & ICOBI & ICOHP (eV) & ICOHP antibonding & ICOBI antibonding \\
\midrule
\endfirsthead
\toprule
Formula & mp-id & Space group & ICOBI & ICOHP (eV) & ICOHP antibonding & ICOBI antibonding \\
\midrule
\endhead
\bottomrule
\endfoot
\bottomrule
\multicolumn{7}{@{}p{0.95\linewidth}}{\footnotesize Black $^*$ = stable experimental compound (thermodynamic phase); red \textcolor{red}{$^*$} = metastable experimental compound; unmarked = theoretical. ICOBI/ICOHP: value of the dominant first-shell species pair (not a mean over every LOBSTER-reported neighbor, see \S\ref{sec:bondtype}).} \\
\endlastfoot
Ba$_{5}$P$_{4}$\textcolor{red}{$^*$} & mp-17566 & Pnma & 0.9040 & -3.4927 & 0.2104 & 0.0000 \\
BaP$_{10}$\textcolor{red}{$^*$} & mp-28035 & Cmc2\_1 & 0.9058 & -4.8297 & 0.0394 & 0.0000 \\
BaS$_{3}$\textcolor{red}{$^*$} & mp-556296 & P2\_12\_12 & 0.9201 & -5.2850 & 0.1084 & 0.0000 \\
Bi$_{2}$O$_{3}$ & mp-674644 & P1 & 0.9332 & -7.5456 & 0.0131 & 0.0000 \\
C(NO)$_{2}$\textcolor{red}{$^*$} & mp-1197051 & P2\_12\_12\_1 & 1.7617 & -16.7609 & 0.0043 & 0.0000 \\
C$_{2}$N$_{2}$O & mp-1214250 & P6\_3/m & 2.4576 & -22.4067 & 0.0003 & 0.0000 \\
CaN$_{6}$ & mp-656097 & C2 & 1.8852 & -17.7987 & 0.0004 & 0.0000 \\
CaN$_{6}$\textcolor{red}{$^*$} & mp-676 & Fddd & 1.8841 & -18.1451 & 0.0003 & 0.0000 \\
CdO$_{2}$\textcolor{red}{$^*$} & mp-2310 & Pa-3 & 0.9738 & -7.1680 & 0.0540 & 0.0000 \\
CdP$_{4}$$^*$ & mp-7904 & P2\_1/c & 0.9223 & -5.1380 & 0.0094 & 0.0000 \\
Cs$_{2}$As$_{3}$$^*$ & mp-15557 & Fmmm & 1.0431 & -4.4248 & 0.0386 & 0.0000 \\
CsN$_{3}$$^*$ & mp-510557 & I4/mcm & 1.8888 & -17.7624 & 0.0000 & 0.0000 \\
CsN$_{3}$ & mp-1225892 & P4/mmm & 1.8983 & -17.5725 & 0.0000 & 0.0000 \\
CsN$_{3}$$^*$ & mp-510557 & I4/mcm & 1.8888 & -17.7624 & 0.0000 & 0.0000 \\
CsN$_{3}$ & mp-1225892 & P4/mmm & 1.8983 & -17.5725 & 0.0000 & 0.0000 \\
CsP$_{7}$$^*$ & mp-1201790 & Pbcm & 0.9343 & -4.9186 & 0.0062 & 0.0000 \\
CsP$_{7}$$^*$ & mp-1201790 & Pbcm & 0.9343 & -4.9186 & 0.0062 & 0.0000 \\
CsP$_{7}$ & mp-1213206 & Pca2\_1 & 0.9336 & -4.9318 & 0.0064 & 0.0000 \\
CsP$_{7}$ & mp-1213206 & Pca2\_1 & 0.9336 & -4.9318 & 0.0064 & 0.0000 \\
H$_{2}$CN$_{5}$ & mp-1212563 & C2/c & 1.2926 & -13.4719 & 0.0040 & 0.0000 \\
H$_{9}$CN$_{9}$\textcolor{red}{$^*$} & mp-1202455 & P-1 & 1.2740 & -12.5502 & 0.0042 & 0.0000 \\
HC & mp-995197 & C2/m & 1.2662 & -10.8409 & 0.0000 & 0.0000 \\
HN & mp-1224501 & P-1 & 0.8490 & -8.0912 & 0.0231 & 0.0000 \\
HN\textcolor{red}{$^*$} & mp-23922 & Pmna & 1.8882 & -19.1499 & 0.0000 & 0.0000 \\
HN\textcolor{red}{$^*$} & mp-720515 & P2\_1/c & 1.5928 & -15.9550 & 0.0000 & 0.0000 \\
KN$_{3}$$^*$ & mp-827 & I4/mcm & 1.8889 & -18.4973 & 0.0000 & 0.0000 \\
KN$_{3}$$^*$ & mp-827 & I4/mcm & 1.8889 & -18.4973 & 0.0000 & 0.0000 \\
KP$_{15}$ & mp-1211471 & P-1 & 0.9260 & -5.0473 & 0.0164 & 0.0000 \\
LiN$_{3}$$^*$ & mp-2659 & C2/m & 1.8523 & -20.2194 & 0.0000 & 0.0000 \\
LiP$_{5}$\textcolor{red}{$^*$} & mp-2412 & Pna2\_1 & 0.9037 & -4.9376 & 0.0026 & 0.0000 \\
LiP$_{5}$\textcolor{red}{$^*$} & mp-2412 & Pna2\_1 & 0.9037 & -4.9376 & 0.0026 & 0.0000 \\
LiP$_{5}$ & mp-32760 & Pna2\_1 & 0.8771 & -4.6811 & 0.0064 & 0.0000 \\
LiP$_{5}$ & mp-32760 & Pna2\_1 & 0.8771 & -4.6811 & 0.0064 & 0.0000 \\
LiS$_{4}$ & mp-1143546 & P1 & 0.9663 & -5.6937 & 0.0304 & 0.0000 \\
MgP$_{4}$$^*$ & mp-384 & P2\_1/c & 0.9196 & -5.0216 & 0.0273 & 0.0000 \\
N$_{2}$O\textcolor{red}{$^*$} & mp-1180473 & Pnma & 2.3033 & -22.3518 & 0.0013 & 0.0000 \\
NaN$_{3}$\textcolor{red}{$^*$} & mp-1066400 & R3m & 1.8853 & -18.8113 & 0.0000 & 0.0000 \\
NaN$_{3}$\textcolor{red}{$^*$} & mp-22003 & R-3m & 1.8872 & -19.1626 & 0.0002 & 0.0000 \\
NaN$_{3}$\textcolor{red}{$^*$} & mp-1066400 & R3m & 1.8853 & -18.8113 & 0.0000 & 0.0000 \\
NaN$_{3}$\textcolor{red}{$^*$} & mp-22003 & R-3m & 1.8872 & -19.1626 & 0.0002 & 0.0000 \\
NaN$_{3}$ & mp-634410 & C2 & 1.8691 & -20.2972 & 0.0000 & 0.0000 \\
NaN$_{3}$ & mp-1065265 & C2/m & 2.7322 & -23.2128 & 0.0006 & 0.0215 \\
NaP$_{15}$$^*$ & mp-1194653 & P-1 & 0.9227 & -5.0620 & 0.0186 & 0.0000 \\
NaP$_{5}$\textcolor{red}{$^*$} & mp-31086 & Pnma & 0.9111 & -4.9558 & 0.0285 & 0.0000 \\
NaS$^*$ & mp-2400 & P6\_3/mmc & 0.9510 & -4.7236 & 0.1279 & 0.0000 \\
PI$_{2}$$^*$ & mp-29443 & P-1 & 0.9586 & -4.8005 & 0.1081 & 0.0000 \\
PNO & mp-675742 & P-1 & 1.1751 & -12.3013 & 0.0406 & 0.0000 \\
PbN$_{6}$\textcolor{red}{$^*$} & mp-620058 & Pna2\_1 & 1.8686 & -19.1852 & 0.0008 & 0.0000 \\
PbN$_{6}$\textcolor{red}{$^*$} & mp-667338 & Pnma & 1.8687 & -19.2227 & 0.0007 & 0.0000 \\
RbN$_{3}$\textcolor{red}{$^*$} & mp-581833 & P4/mmm & 1.8969 & -18.3403 & 0.0000 & 0.0000 \\
RbN$_{3}$ & mp-1219567 & Amm2 & 1.3756 & -9.8572 & 0.0216 & 0.0000 \\
RbP$_{15}$ & mp-1209298 & P-1 & 0.9270 & -5.0385 & 0.0156 & 0.0000 \\
RbS\textcolor{red}{$^*$} & mp-558071 & Immm & 0.9920 & -5.0554 & 0.1964 & 0.0019 \\
Si$_{3}$H & mp-995193 & P-3m1 & 0.9370 & -4.6600 & 0.0000 & 0.0000 \\
SiH$_{2}$\textcolor{red}{$^*$} & mp-1192532 & P2\_1/m & 0.9857 & -4.7720 & 0.0000 & 0.0000 \\
Sr$_{3}$As$_{4}$ & mp-678007 & P1 & 0.8522 & -3.0556 & 0.1161 & 0.0000 \\
Sr$_{3}$P$_{4}$$^*$ & mp-14288 & Fdd2 & 0.8910 & -4.1353 & 0.1224 & 0.0000 \\
Sr$_{3}$P$_{4}$ & mp-682953 & P1 & 0.8889 & -4.0705 & 0.1250 & 0.0000 \\
SrN$_{6}$\textcolor{red}{$^*$} & mp-2131 & Fddd & 1.8857 & -18.4334 & 0.0000 & 0.0000 \\
SrO$_{2}$$^*$ & mp-2697 & I4/mmm & 0.9835 & -7.0573 & 0.1841 & 0.0000 \\
SrO$_{2}$\textcolor{red}{$^*$} & mp-726705 & Pnma & 0.9821 & -6.8700 & 0.2389 & 0.0000 \\
SrO$_{2}$ & mp-2763179 & Pnma & 0.9805 & -6.9738 & 0.2327 & 0.0000 \\
\end{longtable}

}

\subsection{Unclassified Compounds}
\label{sec:si-master-unclassified}
{\footnotesize
\input{si_master_list_unclassified_en.tex}
}

\section{LOBSTER Quality Control}
\label{sec:si-qc}

Every LOBSTER run was screened using the fraction of $k$-points failing
projection orthonormalization (\texttt{lobsterout}) and the maximum
deviation reported in \texttt{bandOverlaps.lobster}, with a red-flag
threshold of 0.1 on the latter (main text, Section~2.7). A systematic sweep
of the full 597-structure dataset under the conservative combined criterion
(100\% $k$-point failure \emph{and} band-overlap deviation $>0.5$
\emph{and} robust $z$-score $>10$ on the resulting $\Delta$(ICOHP)) flagged
a single compound, a copper--gold intermetallic (Cu$_4$Au), for exclusion:

\begin{itemize}
  \item 343/343 $k$-points fail orthonormalization (100\%, not a borderline
  case);
  \item maximum \texttt{bandOverlaps.lobster} deviation of 11.4 (versus the
  0.1 red-flag threshold);
  \item physically implausible ICOHP values (third-neighbor Cu--Cu contacts
  at 4.25--5.04~\AA{} with ICOHP values of $+7.6$ to $+11.6$~eV, an order
  of magnitude too large to be physical);
  \item a $\Delta$(ICOHP)/atom of $-380.5$~eV, $\sim$27$\times$ the
  next-worst case in the rest of the population.
\end{itemize}

The second-worst case identified by the sweep is already 29$\times$ less
extreme than Cu$_4$Au and is plausibly genuine, dense $d$-electron-rich
intermetallic chemistry rather than a projection failure, confirming
Cu$_4$Au as a unique quality-control case rather than a symptom of a wider
problem. Every statistic reported in the main text has been computed on the
remaining 588 (respectively 438--439 case-1-reaction) structures; the
excluded compound's own VASP+LOBSTER output is retained in the data
repository for reference.

\section{Compute Cost}
\label{sec:si-cost}

Table~\ref{tab:si-cost} summarizes VASP and LOBSTER wall-clock time across
the 597 structures. Every job ran on a single SLURM node allocation (16
cores): VASP under MPI parallelism, then LOBSTER under OpenMP parallelism on
the same 16 cores, sequentially within the same job.

\begin{table}[htbp]
\centering
\caption{VASP and LOBSTER compute time (597 structures).}
\label{tab:si-cost}
\begin{tabular}{@{}lrr@{}}
\toprule
 & VASP & LOBSTER \\
\midrule
$n$ & 597 & 596\textsuperscript{a} \\
Mean (s) & 188 & 578 \\
Median (s) & 72 & 157 \\
Max (s) & 5715 (1.6 h) & 22555 (6.3 h) \\
Sequential total & 31.2 h & 95.8 h \\
\bottomrule
\end{tabular}
\vspace{0.15cm}
\\{\footnotesize \textsuperscript{a} one of 597 structures had no usable
LOBSTER step (trivial-cell reference element).}
\end{table}

Combined sequential total (VASP + LOBSTER, 597 structures): 127.0~h ---
between two and three orders of magnitude below what a phonon
(supercell-based) or AIMD screening of the same 597 structures, at the
hierarchy described in the main text's Introduction, would require, and
consistent with the descriptor's intended role as an inexpensive
pre-screen rather than a replacement for that hierarchy.

\section{Case-1 Reaction Descriptor Values}
\label{sec:si-reactions}

Section~\ref{sec:si-full-reactions} reproduces the project's complete
517-reaction case-1 (decomposition-into-elements) population: for each
compound, the balanced reaction (with each elemental reference's space
group), the resulting \texttt{ViabilityLabel}
(\texttt{METASTABLE\_VIABLE}/\texttt{UNSTABLE\_NONEXISTENT}/insufficient
data), and $\Delta$(ICOHP)/atom and $\Delta$(ICOHP\textsubscript{antibonding})/atom
(main text, Eqs.~1--2). Sections~\ref{sec:si-endobondic}
and~\ref{sec:si-exobondic} give the endobondic and exobondic subsets,
respectively, for the 208-compound extension batch of the dataset
(Section~\ref{sec:si-dataset}, Track~2), each reaction annotated with its
$E_\text{hull}$. A condensed summary by stability family is given in
Table~\ref{tab:si-family-breakdown}.

\begin{table}[htbp]
\centering
\caption{Median $\Delta$(ICOHP)/atom and endobondic fraction by stability
class ($n=177$; the subset of the dataset carrying a three-way
experimental-metastable/stable/theoretical-metastable family label,
Section~\ref{sec:si-dataset}).}
\label{tab:si-family-breakdown}
\begin{tabular}{@{}lrrr@{}}
\toprule
Class & $n$ & Median $\Delta$(ICOHP)/atom (eV) & \% endobondic \\
\midrule
Experimental metastable & 59 & $-1.038$ & 28.8\% \\
Stable                  & 59 & $-1.391$ & 22.0\% \\
Theoretical metastable  & 59 & $-1.601$ & 13.6\% \\
\bottomrule
\end{tabular}
\vspace{0.15cm}
\\{\footnotesize Kruskal--Wallis across the three groups: $H=2.56$,
$p=0.278$ (not significant at this subset size), but the ordering matches
the expected physical trend: experimentally metastable compounds have the
highest endobondic fraction, theoretical-only metastable compounds the
lowest.}
\end{table}

\subsection{Complete Case-1 Reaction List}
\label{sec:si-full-reactions}
{\footnotesize
\input{si_full_reactions_en.tex}
}

\subsection{Endobondic Reactions (Extension Batch)}
\label{sec:si-endobondic}
{\footnotesize
\begin{longtable}{@{}lrrr@{}}
\caption{Decomposition-into-elements (case 1) reactions labeled \textbf{endobondic} ($\Delta$ICOHP $\geq 0$, products $-$ reactants) among the 208 compounds of the \texttt{extension} campaign. $^*$ = experimental starting compound (unmarked = theoretical).}\\
\toprule
Reaction (balanced) & mp-id / COD & $E_{hull}$ (eV/at) & $\Delta$ICOHP/at. (eV) \\
\midrule
\endfirsthead
\toprule
Reaction (balanced) & mp-id / COD & $E_{hull}$ (eV/at) & $\Delta$ICOHP/at. (eV) \\
\midrule
\endhead
\bottomrule
\endfoot
0.3333 Be$_{3}$N$_{2}$ $\rightarrow$ Be + 0.3333 N$_{2}$ & mp-1017550 & 0.2081 & 0.3960 \\
BeCl$_{2}$$^*$ $\rightarrow$ Be + Cl$_{2}$ & mp-23267 & 0.0087 & 1.8204 \\
BeF$_{2}$$^*$ $\rightarrow$ Be + F$_{2}$ & mp-15951 & 0.0004 & 3.1514 \\
BeO$^*$ $\rightarrow$ Be + 0.5 O$_{2}$ & mp-7599 & 0.0153 & 2.7864 \\
BeO $\rightarrow$ Be + 0.5 O$_{2}$ & mp-1778 & 0.0070 & 4.4391 \\
0.3333 C$_{3}$N$_{4}$$^*$ $\rightarrow$ C + 0.6667 N$_{2}$ & mp-571653 & 0.4981 & 1.1208 \\
0.3333 C$_{3}$N$_{4}$ $\rightarrow$ C + 0.6667 N$_{2}$ & mp-2852 & 0.5033 & 1.8456 \\
0.5 Cs$_{2}$S$^*$ $\rightarrow$ Cs + 0.5 S & mp-1077814 & 0.0667 & 0.9554 \\
CsCl$^*$ $\rightarrow$ Cs + 0.5 Cl$_{2}$ & mp-573697 & 0.0071 & 0.8830 \\
CsF$^*$ $\rightarrow$ Cs + 0.5 F$_{2}$ & mp-8455 & 0.1088 & 3.8464 \\
CsO$_{2}$$^*$ $\rightarrow$ Cs + O$_{2}$ & mp-1441 & 0.0349 & 0.0666 \\
CsO$_{2}$ $\rightarrow$ Cs + O$_{2}$ & mp-1096936 & 0.0940 & 0.7743 \\
KCl$^*$ $\rightarrow$ K + 0.5 Cl$_{2}$ & mp-23289 & 0.0776 & 0.7616 \\
0.5 Li$_{2}$O$^*$ $\rightarrow$ Li + 0.25 O$_{2}$ & mp-1960 & 0.0000 & 2.3366 \\
0.5 Li$_{2}$O $\rightarrow$ Li + 0.25 O$_{2}$ & mp-755894 & 0.0844 & 1.9481 \\
0.5 Li$_{2}$S$^*$ $\rightarrow$ Li + 0.5 S & mp-1125 & 0.0615 & 1.8505 \\
0.5 Li$_{2}$S $\rightarrow$ Li + 0.5 S & mp-557142 & 0.1750 & 1.5052 \\
0.3333 Li$_{3}$N$^*$ $\rightarrow$ Li + 0.1667 N$_{2}$ & mp-2341 & 0.0039 & 0.6165 \\
LiCl$^*$ $\rightarrow$ Li + 0.5 Cl$_{2}$ & mp-22905 & 0.0243 & 2.8030 \\
LiCl $\rightarrow$ Li + 0.5 Cl$_{2}$ & mp-1185319 & 0.0039 & 2.5005 \\
LiF$^*$ $\rightarrow$ Li + 0.5 F$_{2}$ & mp-1138 & 0.0000 & 4.4673 \\
LiF $\rightarrow$ Li + 0.5 F$_{2}$ & mp-1009009 & 0.2875 & 4.5131 \\
LiP$_{5}$$^*$ $\rightarrow$ Li + 5 P & mp-2412 & 0.0077 & 0.9043 \\
LiP$_{5}$ $\rightarrow$ Li + 5 P & mp-32760 & 0.0963 & 0.8692 \\
MgCl$_{2}$$^*$ $\rightarrow$ Mg + Cl$_{2}$ & mp-23210 & 0.0000 & 0.1568 \\
MgF$_{2}$$^*$ $\rightarrow$ Mg + F$_{2}$ & mp-1072956 & 0.0025 & 0.0015 \\
MgF$_{2}$ $\rightarrow$ Mg + F$_{2}$ & mp-1062842 & 0.2636 & 0.0229 \\
0.5 Rb$_{2}$O$^*$ $\rightarrow$ Rb + 0.25 O$_{2}$ & mp-1394 & 0.0032 & 0.7235 \\
RbCl$^*$ $\rightarrow$ Rb + 0.5 Cl$_{2}$ & mp-23299 & 0.0298 & 1.7924 \\
RbF$^*$ $\rightarrow$ Rb + 0.5 F$_{2}$ & mp-2064 & 0.0516 & 3.9244 \\
RbF $\rightarrow$ Rb + 0.5 F$_{2}$ & mp-11718 & 0.0000 & 1.9497 \\
SiO$_{2}$$^*$ $\rightarrow$ Si + O$_{2}$ & mp-9258 & 0.5770 & 3.3332 \\
SiO$_{2}$ $\rightarrow$ Si + O$_{2}$ & mp-10064 & 1.2779 & 2.1735 \\
SiO$_{2}$ $\rightarrow$ Si + O$_{2}$ & mp-556044 & 1.4574 & 0.0958 \\
SiO$_{2}$ $\rightarrow$ Si + O$_{2}$ & mp-638049 & 1.7448 & 0.3849 \\
\end{longtable}

}

\subsection{Exobondic Reactions (Extension Batch)}
\label{sec:si-exobondic}
{\footnotesize
\input{si_extension_exobondic_en.tex}
}

\section{Additional Correlation Figures}
\label{sec:si-figures}

Figures~\ref{fig:si-reaction-icohp} and~\ref{fig:si-antibonding} reproduce,
for reference, the bond-type-resolved correlation plots for the two
baseline descriptors discussed in the main text (Section~3.2): the fully
integrated reaction $\Delta$(ICOHP) and the raw (non-reaction) near-frontier
antibonding population.

\begin{figure}[htbp]
\centering
\includegraphics[width=0.8\textwidth]{reaction_icohp_vs_formation_energy_per_atom_by_bondtype.png}
\caption{Fully integrated reaction $\Delta$(ICOHP) versus
\texttt{formation\_energy\_per\_atom}, by bonding-type subgroup.}
\label{fig:si-reaction-icohp}
\end{figure}

\begin{figure}[htbp]
\centering
\includegraphics[width=0.8\textwidth]{antibonding_vs_formation_energy_by_bondtype.png}
\caption{Raw (non-reaction) near-frontier antibonding population
(ICOHP-based) versus \texttt{formation\_energy\_per\_atom}, by bonding-type
subgroup.}
\label{fig:si-antibonding}
\end{figure}

\end{document}
